% Compile with XeLaTeX or Tectonic. In Overleaf, choose XeLaTeX.
% Self-contained internship CV: no photograph, external images, or custom class.
\documentclass[11pt,a4paper]{article}
\usepackage[margin=18mm,top=16mm,bottom=17mm,footskip=8mm]{geometry}
\usepackage{fontspec}
\setmainfont{texgyrepagella}[
  Extension=.otf,UprightFont=*-regular,BoldFont=*-bold,
  ItalicFont=*-italic,BoldItalicFont=*-bolditalic]
\setsansfont{texgyreheros}[
  Extension=.otf,UprightFont=*-regular,BoldFont=*-bold,
  ItalicFont=*-italic,BoldItalicFont=*-bolditalic,Scale=MatchLowercase]
\usepackage{xcolor}
\usepackage{tabularx,array}
\usepackage{enumitem}
\usepackage{needspace}
\usepackage{fancyhdr}
\usepackage{hyperref}
\definecolor{accent}{HTML}{183D50}
\definecolor{muted}{HTML}{52616A}
\definecolor{rulecolor}{HTML}{B9C7CF}
\hypersetup{
  colorlinks=true,urlcolor=accent,linkcolor=accent,
  pdftitle={SUN Wei - Internship CV},pdfauthor={SUN Wei},
  pdfsubject={CV for internship applications},
  pdfkeywords={Interaction design, AI applications, frontend development, industrial design}
}
\setlength{\parindent}{0pt}
\setlength{\parskip}{0pt}
\setlength{\tabcolsep}{0pt}
\setlength{\emergencystretch}{1.5em}
\renewcommand{\baselinestretch}{1.035}
\raggedbottom
\pagestyle{fancy}
\fancyhf{}
\renewcommand{\headrulewidth}{0pt}
\fancyfoot[L]{{\sffamily\fontsize{8}{10}\selectfont\color{muted}SUN Wei}}
\fancyfoot[R]{{\sffamily\fontsize{8}{10}\selectfont\color{muted}\thepage}}
\setlist[itemize]{
  leftmargin=1.15em,label=\textbullet,labelsep=0.5em,
  topsep=4pt,itemsep=3pt,parsep=0pt,partopsep=0pt
}

% Reusable formatting commands. Edit the document body below to update content.
\newcommand{\cvsection}[1]{%
  \par\addvspace{11pt}\Needspace{4\baselineskip}%
  {\sffamily\bfseries\fontsize{10.3}{12.4}\selectfont\color{accent}%
    \addfontfeatures{LetterSpace=6}\MakeUppercase{#1}\par}%
  \vspace{3pt}{\color{rulecolor}\hrule height 0.45pt}\vspace{6pt}}
\newcommand{\cventry}[3]{%
  \Needspace{5\baselineskip}%
  \textbf{#1}\hfill{\small\color{muted}#2}\par\vspace{2pt}%
  {\small\color{muted}#3\par}}
\newcommand{\cvskill}[2]{%
  \begin{tabularx}{\linewidth}{@{}>{\raggedright\bfseries\arraybackslash}p{43mm}@{\hspace{4mm}}>{\raggedright\arraybackslash}X@{}}#1&#2\end{tabularx}\par\vspace{5pt}}

\begin{document}

% CONTACT
{\sffamily\bfseries\fontsize{25}{28}\selectfont\color{accent}SUN Wei\par}
\vspace{3pt}
{\fontsize{11}{14}\selectfont Industrial Design undergraduate \textperiodcentered\ Zhejiang University\par}
\vspace{7pt}
{\small
\begin{tabularx}{\linewidth}{@{}X>{\raggedleft\arraybackslash}p{62mm}@{}}
  \textbf{Email:} \href{mailto:1162100953@qq.com}{1162100953@qq.com}
    &\textbf{Phone:} \href{tel:+8618806820562}{+86 188 0682 0562}\\[3pt]
  \textbf{Personal website:} \href{https://crackthedream.github.io/}{crackthedream.github.io}
    &\textbf{WeChat:} Takodachi866
\end{tabularx}}

\cvsection{Education}
\cventry{Zhejiang University}{Sep 2023 -- Jun 2027}{Industrial Design \textperiodcentered\ Hangzhou, China \textperiodcentered\ Expected graduation: June 2027}
College of Computer Science \& Technology\par
\vspace{3pt}
\begin{tabularx}{\linewidth}{@{}X>{\raggedleft\arraybackslash}p{37mm}@{}}
  \textbf{Overall GPA:} 3.85/4.3 (85.97/100)&\textbf{Ranking:} 16/42\\[2pt]
  \textbf{Third-year GPA:} 4.09/4.3 (89.59/100)&
\end{tabularx}

\cvsection{Internship Interests}
User experience, interactive products, and AI applications across design and development.

\cvsection{Research Experience}
\textbf{Speed, Not Density: Oculomotor Responses to Crowd Motion in Passive VR Viewing}\par
\vspace{3pt}
{\small\color{muted}Undergraduate research \textperiodcentered\ First author \textperiodcentered\ Manuscript under review at ACM CHI 2027\par}
\begin{itemize}
  \item Designed a passive-viewing VR study to examine how crowd motion relates to eye movements and discomfort; built 21 Unity stimulus configurations across three crowd behaviors.
  \item Conducted sessions with 42 participants, collecting eye tracking, peripheral physiological signals, and symptom ratings; retained 35 participants and 245 trials after quality screening.
  \item Developed the workflow for multimodal data alignment, cleaning, and eye-movement feature extraction; performed paired nonparametric comparisons and multiple-comparison correction, and wrote the manuscript.
  \item Found higher saccade rates at faster crowd speeds in the primary paired comparison (\textit{n} = 11): 11.55 versus 20.42/min at 3.0 versus 4.5 m/s (BH-adjusted \textit{p} = .0366).
\end{itemize}

\cvsection{Selected Projects}
\cventry{SAVOUR}{May -- Jun 2026}{AR food-information interface \textperiodcentered\ Team project}
\begin{itemize}
  \item Built the recognition-to-overlay software pipeline for an AR food-information prototype, connecting detected food regions to nutrition feedback and taking responsibility for software integration and debugging.
  \item Integrated GLEE region detection with Qwen visual recognition; revised category correction, repeated detection, short-term tracking, and label expiration to address ambiguous targets and drifting or stale overlays.
\end{itemize}
\vspace{4pt}
\cventry{Momentum}{Apr 2026}{AI study support and self-management \textperiodcentered\ Team project}
\begin{itemize}
  \item Built the React/Vite mobile web frontend and connected FastAPI services for task planning, priorities, deadlines, and focused study sessions; a teammate provided the original visual design.
  \item Integrated LLM tool calls with stored task data and connected speech recognition and synthesis for contextual dialogue during focus sessions; used mode-specific prompts to distinguish study support from task-editing requests.
\end{itemize}

\newpage
\cventry{EMS Feedback for Multi Drone Search Simulation}{Jan 2026}{Simulation and physical feedback \textperiodcentered\ Team project}
\begin{itemize}
  \item Implemented communication that turns target events in a Unity drone-search simulation into forearm electrical muscle stimulation (EMS), encoding target direction and distance as physical feedback.
  \item Connected Unity events to a Python device service through TCP/JSON and sent commands over Bluetooth Low Energy; contributed to feedback mapping and demonstrated the integrated pipeline in the course project.
\end{itemize}

\cvsection{Additional Design Experience}
\cventry{FOCUS DRIVE}{Dec 2025}{Driver-service interface \textperiodcentered\ Team project}
\vspace{4pt}
Designed feature flows and six groups of high-fidelity screens for a driver-service app, covering order handling, navigation, and safety-related information.\par
\vspace{9pt}
\cventry{X-Band}{Jul 2025}{Wearable interaction concept \textperiodcentered\ Individual project}
\vspace{4pt}
Designed a wearable concept for people with limb loss to log phantom-limb discomfort and access professionally configured stimulation; developed one-handed controls, user confirmation, and an explicit stop mechanism.

\cvsection{Honors}
\textbf{Academic Excellence Exemplar}\hfill{\small\color{muted}2026}\par

\cvsection{Skills}
\cvskill{Frontend and AI}{React/Vite frontend development; FastAPI service integration; LLM tool calls and computer vision application integration.}
\cvskill{Interactive systems}{Unity prototyping; software-hardware integration; TCP/JSON and Bluetooth Low Energy device communication.}
\cvskill{Interaction design}{Interaction flows and Figma prototyping; user studies and quantitative analysis with Python; Blender 3D modeling and visualization.}

\cvsection{Languages}
\textbf{Mandarin Chinese:} Native.\quad \textbf{English:} CET-6, 609/710.\par
\vspace{3pt}
\textbf{TOEFL iBT:} 5.0/6 (100/120; CEFR C1), 30 September 2026.\par
Reading 5.0, Listening 5.0, Writing 5.5, Speaking 4.5 (each out of 6).

\end{document}
